\begin{thebibliography}{99}
\footnotesize\itemsep=0pt

\bibitem{p-C14}
Cuntz M., Frieze patterns as root posets and affine triangulations,
 \href{https://doi.org/10.1016/j.ejc.2014.06.005}{\textit{European~J.~Combin.}} \textbf{42} (2014), 167--178,
 \href{https://arxiv.org/abs/1307.7986}{arXiv:1307.7986}.

\bibitem{p-CH09b}
Cuntz M., Heckenberger I., Weyl groupoids of rank two and continued fractions,
 \href{https://doi.org/10.2140/ant.2009.3.317}{\textit{Algebra Number Theory}} \textbf{3} (2009), 317--340,
 \href{https://arxiv.org/abs/0807.0124}{arXiv:0807.0124}.

\bibitem{p-CH09a}
Cuntz M., Heckenberger I., Weyl groupoids with at most three objects,
 \href{https://doi.org/10.1016/j.jpaa.2008.11.009}{\textit{J.~Pure Appl. Algebra}} \textbf{213} (2009), 1112--1128,
 \href{https://arxiv.org/abs/0805.1810}{arXiv:0805.1810}.

\bibitem{p-CH10}
Cuntz M., Heckenberger I., Finite {W}eyl groupoids, \href{https://doi.org/10.1515/crelle-2013-0033}{\textit{J.~Reine Angew.
 Math.}} \textbf{702} (2015), 77--108, \href{https://arxiv.org/abs/1008.5291}{arXiv:1008.5291}.

\bibitem{p-CL-17}
Cuntz M., Lentner S., A simplicial complex of {N}ichols algebras,
 \href{https://doi.org/10.1007/s00209-016-1711-0}{\textit{Math.~Z.}} \textbf{285} (2017), 647--683, \href{https://arxiv.org/abs/1503.08117}{arXiv:1503.08117}.

\bibitem{Eilenberg1950}
Eilenberg S., MacLane S., Cohomology theory of {A}belian groups and homotopy
 theory.~{I}, \href{https://doi.org/10.1073/pnas.36.8.443}{\textit{Proc. Nat. Acad. Sci. USA}} \textbf{36} (1950), 443--447.

\bibitem{Eilenberg1950a}
Eilenberg S., MacLane S., Cohomology theory of {A}belian groups and homotopy
 theory.~{II}, \href{https://doi.org/10.1073/pnas.36.11.657}{\textit{Proc. Nat. Acad. Sci. USA}} \textbf{36} (1950),
 657--663.

\bibitem{b-EGNO-15}
Etingof P., Gelaki S., Nikshych D., Ostrik V., Tensor categories, \textit{Math.
 Surveys Monogr.}, Vol.~205, \href{https://doi.org/10.1090/surv/205}{Amer. Math. Soc.}, Providence, RI, 2015.

\bibitem{p-H-06}
Heckenberger I., The {W}eyl groupoid of a {N}ichols algebra of diagonal type,
 \href{https://doi.org/10.1007/s00222-005-0474-8}{\textit{Invent. Math.}} \textbf{164} (2006), 175--188,
 \href{https://arxiv.org/abs/math.QA/0411477}{arXiv:math.QA/0411477}.

\bibitem{p-H-09}
Heckenberger I., Classification of arithmetic root systems, \href{https://doi.org/10.1016/j.aim.2008.08.005}{\textit{Adv. Math.}}
 \textbf{220} (2009), 59--124, \href{https://arxiv.org/abs/math.QA/0605795}{arXiv:math.QA/0605795}.

\bibitem{p-J-77}
Joseph A., A generalization of the {G}elfand--{K}irillov conjecture,
 \href{https://doi.org/10.2307/2374020}{\textit{Amer.~J.~Math.}} \textbf{99} (1977), 1151--1165.

\bibitem{b-Kac90}
Kac V.G., Infinite-dimensional {L}ie algebras, 3rd~ed., \href{https://doi.org/10.1017/CBO9780511626234}{Cambridge University
 Press}, Cambridge, 1990.

\bibitem{MR4498166}
Laugwitz R., Walton C., Constructing non-semisimple modular categories with
 relative monoidal centers, \href{https://doi.org/10.1093/imrn/rnab097}{\textit{Int. Math. Res. Not.}} \textbf{2022}
 (2022), 15826--15868, \href{https://arxiv.org/abs/2010.11872}{arXiv:2010.11872}.

\bibitem{MacLane1952}
MacLane S., Cohomology theory of {A}belian groups, in Proceedings of the
 {I}nternational {C}ongress of {M}athematicians, {C}ambridge, {M}ass., 1950,
 Vol.~2, Amer. Math. Soc., Providence, R.I., 1952, 8--14.

\bibitem{MR1632802}
Rosso M., Quantum groups and quantum shuffles, \href{https://doi.org/10.1007/s002220050249}{\textit{Invent. Math.}}
 \textbf{133} (1998), 399--416.

\end{thebibliography}
